\documentclass[11pt]{article}
\usepackage[margin=1in]{geometry}
\usepackage{amsmath,amssymb}

\begin{document}

\begin{center}
  \textbf{\large Derivation of the deformed generating function
  $H^{(c)}(t)=\dfrac{e^{ct}}{1-t^2}$}
\end{center}

\subsection*{Question}

Take the sequence with $h_0 = 1$ and $h_r = 2$ for $r \ge 1$.  Its generating
function is $H(t) = \frac{1+t}{1-t}$, which forces the paired sequence
$j = (2,0,2,0,\dots)$.  The $c$-deformation prepends $c$ to $j$.  Working
through equation~(D), show that the deformed generating function is
\[
  H^{(c)}(t) = \frac{e^{ct}}{1-t^2}.
\]

\subsection*{The relation (D)}

Everything rests on equation~(D) written as a logarithmic derivative.  With
$H(t) = \sum_{n\ge 0} h_n t^n$ and $J(t) = \sum_{r\ge 1} j_r t^r$, the
following three statements are equivalent:
\[
  J(t) = t\,\frac{d}{dt}\log H(t),
  \qquad
  t H'(t) = H(t)\,J(t),
  \qquad
  n h_n = \sum_{r=1}^{n} j_r\, h_{n-r}\ \ (n \ge 1).
\]
The middle identity follows by matching the coefficient of $t^n$ on each
side of the first; the third is that coefficient identity written out.

\subsection*{Step 1: the base $J$}

With $h_0 = 1$ and $h_r = 2$ for $r \ge 1$,
\[
  H(t) = 1 + 2\sum_{r\ge 1} t^r = 1 + \frac{2t}{1-t} = \frac{1+t}{1-t}.
\]
Hence
\[
  \frac{d}{dt}\log H(t)
    = \frac{d}{dt}\bigl[\log(1+t) - \log(1-t)\bigr]
    = \frac{1}{1+t} + \frac{1}{1-t}
    = \frac{2}{1-t^2},
\]
and therefore
\[
  J(t) = t\cdot\frac{2}{1-t^2}
       = \frac{2t}{1-t^2}
       = 2t + 2t^3 + 2t^5 + \cdots,
\]
so that $j = (2,0,2,0,\dots)$, i.e. $j_r = 2$ for odd $r$ and $0$ for even $r$.

\subsection*{Step 2: prepending $c$ to a generating function}

The deformation sends a sequence $S = \{x_1, x_2, \dots\}$ to
$S^{(c)} = \{c, x_1, x_2, \dots\}$: index $1$ becomes $c$, and each old entry
$x_r$ moves to index $r+1$.  If $X(t) = \sum_{r\ge 1} x_r t^r$, then
\[
  X^{(c)}(t)
    = c\,t + \sum_{r\ge 2} x_{r-1} t^{r}
    = c\,t + t\sum_{s\ge 1} x_s t^{s}
    = c\,t + t\,X(t).
\]
Applying this to $j$,
\[
  J^{(c)}(t) = c\,t + t\,J(t)
             = c\,t + \frac{2t^2}{1-t^2}.
\]

\subsection*{Step 3: recover $H^{(c)}$ from $J^{(c)}$}

The deformed sequence $h^{(c)}$ is the one paired with $j^{(c)}$ through the
same relation~(D), so
\[
  \frac{d}{dt}\log H^{(c)}(t) = \frac{J^{(c)}(t)}{t} = c + \frac{2t}{1-t^2}.
\]
Integrating, and using $\displaystyle\int \frac{2t}{1-t^2}\,dt = -\log(1-t^2)$
(substitute $u = 1-t^2$, $du = -2t\,dt$),
\[
  \log H^{(c)}(t) = c\,t - \log(1-t^2) + C.
\]
The constant $C$ is fixed by the normalization $h^{(c)}_0 = 1$, i.e.
$H^{(c)}(0) = 1$, hence $\log H^{(c)}(0) = 0$.  Evaluating the right-hand side
at $t = 0$ gives $C = 0$.  Exponentiating,
\[
  \boxed{\,H^{(c)}(t) = \frac{e^{ct}}{1-t^2}\,.}
\]

\subsection*{Remark on the integration constant}

The antiderivative $\int \frac{2t}{1-t^2}\,dt$ is determined only up to the
branch choice $-\log(1-t^2)$ versus $-\log(t^2-1)$, which differ by a
constant and produce $e^{ct}/(1-t^2)$ versus $e^{ct}/(t^2-1)$.  These differ
by an overall sign.  The boundary condition $H^{(c)}(0) = 1$ selects the
former: at $t = 0$ it equals $+1$, whereas $e^{ct}/(t^2-1)$ equals $-1$.  A
symbolic integrator that fixes its own constant of integration may return the
second form; the physically correct branch is $e^{ct}/(1-t^2)$.

\subsection*{Consistency check}

One verifies directly that $H^{(c)}(t) = e^{ct}/(1-t^2)$ reproduces
$J^{(c)}(t)$ through relation~(D):
\[
  t\,\frac{d}{dt}\log H^{(c)}(t)
    = t\,\frac{d}{dt}\bigl[c\,t - \log(1-t^2)\bigr]
    = t\left(c + \frac{2t}{1-t^2}\right)
    = c\,t + \frac{2t^2}{1-t^2}
    = J^{(c)}(t),
\]
as required.

\end{document}
